% !TEX root = ../Main.tex
\chapter{回归教材}

立体几何证明题靠的不是凭空想象，而是把图里的条件翻译成判定定理用得上的形式。所以第一步是回归教材：把基本概念、公理、推论和各种位置关系先理顺，做题时才知道每一步在调用哪条定理。

\section{基本概念与表示方法}

\state{基本概念}{
    点；线（点构成的集合）；面（点或线构成的集合）。
}

\state{基本表示方法}{
    点 $A$ 在直线 $l$ 上：$A\in l$；\qquad 在平面 $\alpha$ 内：$A\in\alpha$。

    直线 $l$ 在平面 $\alpha$ 内：$l\subset\alpha$。
}

\section{公理与推论}

\state{公理}{
    \begin{enumerate}
        \item 不共线的三点确定一个平面；
        \item 若一条直线上的两个点在一个平面内，则此直线在这个平面内；
        \item 若两个不重合的平面有一个公共点，则它们有且只有一条过该点的公共直线。
    \end{enumerate}
}

由公理推出三条常用推论——平时"确定一个平面"最常用的依据：

\state{推论}{
    \begin{enumerate}
        \item 经过一条直线和这条直线外一点，有且只有一个平面；
        \item 经过两条相交直线，有且只有一个平面；
        \item 经过两条平行直线，有且只有一个平面。
    \end{enumerate}
}

\section{三类位置关系}

把空间里直线、平面的相对位置分类清楚，是后面判定平行、垂直的前提。

\textbf{两条直线}的位置关系：

\begin{center}
\begin{tikzpicture}[>=Stealth, every node/.style={font=\small}]
    \node (a) at (0,0) {两条直线};
    \node (b1) at (3.2,0.7) {异面直线};
    \node (b2) at (3.2,-0.7) {共面直线};
    \node (c1) at (6.4,-0.2) {相交直线};
    \node (c2) at (6.4,-1.2) {平行直线};
    \draw (a) -- (b1); \draw (a) -- (b2);
    \draw (b2) -- (c1); \draw (b2) -- (c2);
\end{tikzpicture}
\end{center}

\textbf{直线与平面}的位置关系：

\begin{center}
\begin{tikzpicture}[>=Stealth, every node/.style={font=\small}]
    \node (a) at (0,0) {直线与平面};
    \node (b1) at (3.4,0.7) {在平面内};
    \node (b2) at (3.4,-0.7) {在平面外};
    \node (c1) at (6.6,-0.2) {相交};
    \node (c2) at (6.6,-1.2) {平行};
    \draw (a) -- (b1); \draw (a) -- (b2);
    \draw (b2) -- (c1); \draw (b2) -- (c2);
\end{tikzpicture}
\end{center}

\textbf{两个平面}的位置关系：

\begin{center}
\begin{tikzpicture}[>=Stealth, every node/.style={font=\small}]
    \node (a) at (0,0) {两个平面};
    \node (b1) at (3.2,0.55) {平行};
    \node (b2) at (3.2,-0.55) {相交};
    \draw (a) -- (b1); \draw (a) -- (b2);
\end{tikzpicture}
\end{center}

后面两章就围绕"平行"和"相交（垂直）"这两类关系，把线线、线面、面面三个层次的判定与性质串起来。
